\Needspace{18\baselineskip}
\section{The abstract specimen interface}

\begin{definition}[Focus modes]\label{term:focus}
The type of modes has exactly three constructors:
$\mathsf{search}$, $\mathsf{backtrack}$, and $\mathsf{rewrite}$.
\end{definition}
\lean{LocalReasoning}{15}{17}{Focus}
They represent the paper's searching, backtracking, and triggering focuses.
The derived decidable equality allows their distinctness to be checked by
finite case analysis. This datatype records only a mode, not a graph
position or a focus edge.

\begin{definition}[Specimen system]\label{term:specimen}\label{def:specimen}
A specimen system $G$ over $M$ consists of a type $\Sigma$ of specimens,
maps $\ell,r:\Sigma\to X$, a mode map
$\mode:\Sigma\to\mathsf{Focus}$, a predicate $\ent$ on specimens,
and a predicate $\rt$ on states. Its laws are:
\begin{enumerate}
\item \emph{Rootedness preservation:} $\rt(s)$ and $s\to t$ imply $\rt(t)$.
\item \emph{Final-state agreement:} for every $c\in\Sigma$,
  $\fin(\ell(c))\iff\fin(r(c))$.
\item \emph{Context-only traversal:} suppose both endpoints of $c$ are
  rooted and $c$ is either backtracking or searching without entering.
  If either endpoint has an outgoing transition, there is a specimen $d$
  whose focus is remote and such that
  \[
    \ell(c)\to\ell(d),\qquad r(c)\to r(d).
  \]
\end{enumerate}
The terms ``rooted,'' ``entering,'' and ``remote'' receive their precise
abstract meanings from these data and the definitions immediately below.
\end{definition}
\lean{LocalReasoning}{24}{38}{SpecimenSystem}
The corresponding fields are \texttt{Specimen}, \texttt{left},
\texttt{right}, \texttt{mode}, \texttt{entering}, \texttt{rooted},
\texttt{rooted\_step}, \texttt{same\_final}, and \texttt{context\_step}.
The last three are hypotheses needed to construct an instance, not results
proved here from a graph model. The intended specimen consists of a
multi-hole focussed context and two pointwise related vectors of fragments.
That intended representation is not a field of this abstract structure.

In particular, the interface does not identify modes with a focus field on
states, define rootedness as reachability from an initial state, or define
entering by incidence with a hole. The context-traversal law captures the
consequence of graph locality and output-closure that the simulation proof
uses. Its availability premise is symmetric between the two endpoints,
which will be essential in the converse proof.

\begin{definition}[Remote, settled, and lifted specimens]\label{term:lifting}
A specimen $c$ is remote if
\[
 \remote(c)\iff \mode(c)\ne\mathsf{rewrite}\ \lor\ \neg\ent(c).
\]
It is settled if
\[
 \settled(c)\iff \mode(c)=\mathsf{backtrack}\ \lor\
   (\mode(c)=\mathsf{search}\land\neg\ent(c)).
\]
Its lifted state relation is
\[
 \Lft(s,t)\iff\exists c\in\Sigma\;
  \remote(c)\land\rt(\ell(c))\land\rt(r(c))
  \land s=\ell(c)\land t=r(c).
\]
\end{definition}
\lean{LocalReasoning}{45}{45}{SpecimenSystem.Remote}
\lean{LocalReasoning}{47}{48}{SpecimenSystem.Settled}
\lean{LocalReasoning}{50}{52}{SpecimenSystem.Lifting}
Remote excludes only an entering rewrite focus. It allows an entering
search focus. Settled is stronger: it excludes every rewrite focus and
also excludes an entering search focus. The two equalities in the lifting
definition expose a particular specimen as a witness for the endpoints.

\begin{lemma}[Settled implies remote]
If $\settled(c)$, then $\remote(c)$.
\end{lemma}
\lean{LocalReasoning}{54}{56}{SpecimenSystem.settled_remote}
\begin{proof}
A backtracking or searching mode is distinct from the rewrite mode.
Therefore the first disjunct in the remote predicate holds in either case.
\end{proof}

\begin{lemma}[Rootedness along a finite execution]
If $\rt(s)$ and $s\steps{n}t$, then $\rt(t)$.
\end{lemma}
\lean{LocalReasoning}{58}{62}{SpecimenSystem.rooted_steps}
\begin{proof}
Induct on the execution. At length zero nothing changes. For a positive
execution, the one-step rootedness law makes its first successor rooted;
the induction hypothesis propagates this property through the remainder.
\end{proof}

\begin{definition}[Abstract non-rewrite quasi-specimen]\label{term:quasi}
For relations $A,B$ on states, define $\quasi{A}{B}(s,t)$ to mean that
\emph{if} $s$ and $t$ are rooted, then there is a specimen $c$ with
\[
 \begin{gathered}
 \mode(c)\ne\mathsf{rewrite},\qquad
 \rt(\ell(c)),\quad\rt(r(c)),\\
 A(s,\ell(c)),\qquad B(r(c),t).
 \end{gathered}
\]
\end{definition}
\lean{LocalReasoning}{67}{70}{SpecimenSystem.Quasi}
The outer implications matter: no witness is required if either endpoint
is unrooted. Subsequent proofs establish both rootedness hypotheses before
using a quasi-specimen. Unlike the paper's full definition, this abstract
consequence does not compare focus labels of $s,t$ with that of $c$;
there is no state-focus map in the interface. It specifically packages the
non-rewrite situation used by safety and robustness.

\begin{lemma}[A rooted quasi-specimen supplies an up-to witness]
If $\quasi{\preceq_P}{\preceq_S}(s,t)$ and both $s,t$ are rooted, then
$\upto{P}{S}{\Lft}(s,t)$.
\end{lemma}
\lean{LocalReasoning}{72}{77}{SpecimenSystem.quasi_upTo}
\begin{proof}
Obtain its witness $c$. Its non-rewrite mode makes it remote. Together
with rootedness of its endpoints, this gives $\Lft(\ell(c),r(c))$.
Use $\ell(c)$ and $r(c)$ as the intermediate states in the up-to
definition, with the two refinements furnished by the quasi-specimen.
\end{proof}

\Needspace{24\baselineskip}
\begin{definition}[Input-safety]\label{term:safety}
The predicate $\safe{Q}{P}$ requires the following for every specimen $c$
whose endpoints are rooted, whose mode is search, and which is entering.
At least one of these three cases must hold:
\begin{enumerate}[label=\textup{(\Roman*)}]
\item Both endpoints have finite executions to stuck states. Their counts
  may be unrelated, and either count may be zero.
\item For some settled specimen $d$ and $k,l\in\mathbb N$,
  \[
    \ell(c)\steps{1+k}\ell(d),\quad r(c)\steps{l}r(d),\quad
    (1+k)\,Q\,l.
  \]
\item For some states $a,b$ and $k,l\in\mathbb N$,
  \[
    \ell(c)\steps{1+k}a,\quad r(c)\steps{1+l}b,\quad
    (1+k)\,Q\,(1+l),
  \]
  and $\quasi{\simeq_P}{\simeq_P}(a,b)$ holds.
\end{enumerate}
\end{definition}
\lean{LocalReasoning}{81}{91}{SpecimenSystem.InputSafe}
Case~II requires a positive left execution but permits zero right steps.
Case~III requires positive executions on both sides and auxiliary
\emph{equivalences}. Retaining both directions of those equivalences is
necessary for reusing safety in the converse theorem.

\begin{definition}[Robustness]\label{term:robustness}
The predicate $\robust{Q}{P}{S}$ requires that for every specimen $c$
with rooted endpoints, rewrite mode, and $\neg\ent(c)$, each transition
$\ell(c)\to s'$ satisfies one of:
\begin{enumerate}[label=\textup{(\roman*)}]
\item $s'$ has a finite execution to a stuck state;
\item for some $k,l,a,b$,
  \[
   s'\steps{k}a,\quad r(c)\steps{l}b,\quad (1+k)\,Q\,l,
  \]
  and $\quasi{\preceq_P}{\preceq_S}(a,b)$ holds.
\end{enumerate}
\end{definition}
\lean{LocalReasoning}{95}{101}{SpecimenSystem.Robust}
This is a rooted, remote restriction of the consequences used from the
paper's robustness condition. The paper's alternative that an endpoint is
unrooted is contradictory under the present premises. The abstract
machine does not classify its transition relation into search and rewrite
rules; it is the specimen's mode that determines where this predicate is
invoked. Establishing the correspondence with concrete rewrite transitions
is part of the missing hypernet instance.
